% Section 3 -- Problem Formulation (P2)
% Needs: \usepackage{amsmath, amssymb, booktabs}
% Every number here comes from src/problem.py or results/lambda_sweep.csv.

\section{Problem Formulation}
\label{sec:formulation}

This section turns the scheduling problem into the cost function that QAOA minimises.
Table~\ref{tab:notation} lists the notation used throughout the report.

\begin{table}[h]
  \centering
  \caption{Notation.}
  \label{tab:notation}
  \begin{tabular}{@{}ll@{}}
    \toprule
    Symbol & Meaning \\
    \midrule
    $n$ & number of jobs (= number of qubits) \\
    $\tilde w_i$ & runtime of job $i$ in raw units (integer) \\
    $w_i = \tilde w_i / \max_k \tilde w_k$ & normalised runtime, $w_i \in (0, 1]$ \\
    $z_i \in \{+1, -1\}$ & assignment of job $i$: $+1$ = node A, $-1$ = node B \\
    $x_i \in \{0, 1\}$ & measured bit for qubit $i$, with $z_i = 1 - 2x_i$ \\
    $E$ & conflict (anti-affinity) graph: pairs that should not share a node \\
    $D = \sum_i w_i z_i$ & load difference between node A and node B \\
    $V(z)$ & number of violated conflicts (edges with both jobs on one node) \\
    $\lambda$ & conflict penalty weight \\
    $J_{ij}$ & Ising coupling between qubits $i$ and $j$ \\
    $s = \max_{i<j} |J_{ij}|$ & Hamiltonian scale factor \\
    $C(z)$, $H_C$ & classical cost function and its cost Hamiltonian \\
    \bottomrule
  \end{tabular}
\end{table}

\subsection{Scheduling model}

We place $n$ jobs on two identical compute nodes, A and B.
Job $i$ has runtime $\tilde w_i$, and some pairs of jobs conflict: for example, two memory-bound jobs that slow each other down when co-located.
The goal is to finish all jobs as early as possible while keeping conflicting jobs on different nodes.

Each job is described by one binary variable $z_i$.
With total work $W = \sum_i w_i$, the loads on the two nodes are $W_A = (W + D)/2$ and $W_B = (W - D)/2$, so the time at which the busier node finishes (the \emph{makespan}) is
\begin{equation}
  \text{makespan} = \max(W_A, W_B) = \frac{W + |D|}{2}.
  \label{eq:makespan}
\end{equation}
Since $W$ does not depend on the assignment, minimising the makespan is the same as minimising $D^2$.
Without conflicts this is the number partitioning problem, the first example in Lucas's catalogue of Ising formulations~\cite{lucas2014}.
Its decision version is NP-complete, although it is only weakly NP-hard: pseudo-polynomial dynamic programming and the Karmarkar--Karp differencing heuristic~\cite{karmarkar1982} solve typical instances well.
We compare against these methods in Section~\ref{sec:results}.

\subsection{Instance generation}
\label{sec:instances}

Instances are generated from a single integer seed, so every method in this report is run on exactly the same problems.

\paragraph{Runtimes.}
Real HPC job runtimes are heavy-tailed: most jobs are short and a few are very long.
We draw $\tilde w_i = \max\!\big(1, \operatorname{round}(e^{X_i})\big)$ with $X_i \sim \mathcal{N}(2, 1)$, a lognormal distribution.
Over 3{,}200 sampled jobs the median runtime is 8, the mean is 12, the 95th percentile is 37 and the largest is 205.

\paragraph{Conflicts.}
Our original plan drew the conflict graph as an Erd\H{o}s--R\'enyi graph, keeping each pair with probability $q = 0.2$.
With two nodes, however, all conflicts can be satisfied only if the graph is bipartite.
A single triangle of mutually conflicting jobs makes a conflict-free schedule impossible, and a random graph with $q = 0.2$ contains $\binom{n}{3} q^3$ triangles on average: about 1 at $n = 10$ and 9 at $n = 20$.
We therefore use a \emph{planted bipartition}: the jobs are split at random into two hidden halves of sizes $\lfloor n/2 \rfloor$ and $\lceil n/2 \rceil$, and each pair with one job in each half becomes a conflict with probability $q = 0.2$.
The hidden split is guaranteed to have zero violations, so every instance has at least one feasible schedule.
The expected number of conflicts is $q \lfloor n/2 \rfloor \lceil n/2 \rceil$, for example 5 at $n = 10$ and 12.8 at $n = 16$.
The runtimes are drawn independently of the hidden split, so it is not a balanced schedule; the solver still has to trade balance against conflicts.

\subsection{Cost function and cost Hamiltonian}

The cost function combines the balance objective with a penalty for conflicts:
\begin{equation}
  C(z) = \Big( \sum_{i=1}^{n} w_i z_i \Big)^2 + \lambda \sum_{(i,j) \in E} z_i z_j .
  \label{eq:cost}
\end{equation}
For a conflict edge, $z_i z_j = +1$ if both jobs are on the same node (a violation) and $-1$ otherwise.
Summing over the edges gives $\sum_{(i,j)\in E} z_i z_j = 2V(z) - |E|$, so the second term equals $2\lambda V(z)$ up to a constant: each violated conflict adds $2\lambda$ to the cost.
This term is an antiferromagnetic Ising model on $E$, the same structure as Max-Cut~\cite{farhi2014}.

Expanding the square and using $z_i^2 = 1$ puts the cost in Ising form, with only pairwise terms:
\begin{equation}
  C(z) = \sum_{i} w_i^2 + \sum_{i<j} J_{ij}\, z_i z_j,
  \qquad
  J_{ij} = 2 w_i w_j + \lambda \, \mathbb{1}[(i,j) \in E].
  \label{eq:ising}
\end{equation}
The balance term couples every pair of jobs, so the Ising graph is complete, with $n(n-1)/2$ couplings; conflicts strengthen the couplings on their edges.

To obtain the cost Hamiltonian we replace each $z_i$ by the Pauli operator $Z_i$ acting on qubit $i$.
Because $Z\lvert 0\rangle = +\lvert 0\rangle$ and $Z\lvert 1\rangle = -\lvert 1\rangle$, measuring bit $x_i = 0$ corresponds to $z_i = +1$ (node A) and $x_i = 1$ to $z_i = -1$ (node B).
The Hamiltonian
\begin{equation}
  H_C = \frac{1}{s} \Big( \sum_i w_i^2 \, I + \sum_{i<j} J_{ij}\, Z_i Z_j \Big)
  \label{eq:hamiltonian}
\end{equation}
is diagonal in the computational basis, and its eigenvalue on $\lvert x \rangle$ is exactly $C(z)/s$.
Finding the lowest-energy basis state of $H_C$ is therefore the same as finding the best schedule.
In the QAOA circuit each term $J_{ij} Z_i Z_j$ becomes one $ZZ$ phase gate per layer, so a layer contains $n(n-1)/2$ two-qubit gates.
The constant term does not change which schedule is best and adds no gates.

\subsection{Normalisation}
\label{sec:normalisation}

We rescale twice before building the circuit.

\begin{enumerate}
  \item \textbf{Runtimes} are divided by the largest runtime, $w_i = \tilde w_i / \max_k \tilde w_k$, so $w_i \in (0,1]$.
  Raw runtimes reach the hundreds, which would make $J_{ij}$ of order $10^4$.
  \item \textbf{The Hamiltonian} is divided by $s = \max_{i<j} |J_{ij}|$, so that its largest coupling is exactly 1.
\end{enumerate}

Neither step changes which schedule is optimal; both change only the scale of the energies.
The scale matters for QAOA because the cost layer applies $e^{-i\gamma H_C}$, which rotates each term by an angle $\gamma J_{ij}$.
With large couplings, a tiny change in $\gamma$ produces a large change in phase, and the energy landscape over $\gamma$ oscillates so quickly that the classical optimiser cannot follow it.
Fixing the largest coupling at 1 gives every instance and every problem size the same $\gamma$ scale, so angle bounds, initial guesses and transferred angles (Section~\ref{sec:method}) mean the same thing across instances.
Energies measured on the circuit are converted back to costs with $C = s\,\langle H_C \rangle$ before any metric is computed.

One side effect of heavy-tailed runtimes is that the couplings between short jobs are very small: at $n = 16$ the median instance has a smallest coupling of about $0.002$ after scaling.
Schedules that differ only in where they place short jobs therefore differ only slightly in energy, which makes them hard to tell apart at low circuit depth.
We return to this in Section~\ref{sec:discussion}.

\subsection{Mirror symmetry and degenerate optima}
\label{sec:symmetry}

Swapping the labels of the two nodes, $z \mapsto -z$, never changes the cost, because Eq.~\eqref{eq:ising} contains only products $z_i z_j$.
Every schedule therefore has a mirror image with the same cost, and the optimum is at least two-fold degenerate.
In operator terms, $H_C$ commutes with the global flip $X^{\otimes n}$.
The QAOA mixer $\sum_i X_i$ and the initial state $\lvert + \rangle^{\otimes n}$ are also symmetric under this flip, so the QAOA state gives every bitstring and its mirror exactly the same probability.
This gives a simple check on the implementation: sampled counts should be mirror-symmetric up to shot noise.

The degeneracy can be larger than two, because integer runtimes often allow several different balanced splits.
At $\lambda = 0.5$, between 20\,\% and 45\,\% of our instances (depending on $n$) have more than two optimal bitstrings.
We therefore define the success probability $P(\text{opt})$ as the total probability of \emph{all} optimal bitstrings found by brute force, not just one mirror pair.

The symmetry could be removed by fixing $z_0 = +1$ and using $n - 1$ qubits.
We keep one qubit per job so that bitstrings map directly to schedules and qubit counts match job counts in the scaling results.

\subsection{Choosing the penalty weight $\lambda$}
\label{sec:lambda}

The penalty weight sets the exchange rate between balance and conflicts.
If $\lambda$ is too small, the cheapest schedule may accept a conflict to gain a little balance.
If $\lambda$ is too large, the conflict couplings dominate $s$, and after scaling the balance couplings become small relative to the largest coupling, which makes the balance part of the landscape harder for QAOA to resolve.
We therefore choose the \emph{smallest} $\lambda$ that reliably removes conflicts.

We swept $\lambda \in \{0.1, 0.5, 1, 2, 5\}$ on 20 instances at each size $n \in \{6, 8, \dots, 16\}$, solving each instance exactly by brute force.
The rule was: choose the smallest $\lambda$ for which the optimal schedule has zero violations on at least 90\,\% of instances at every size.
Table~\ref{tab:lambda} shows the result.

\begin{table}[h]
  \centering
  \caption{Fraction of instances whose optimal schedule has zero conflict violations (20 instances per cell).}
  \label{tab:lambda}
  \begin{tabular}{@{}lcccccc@{}}
    \toprule
    $\lambda$ & $n=6$ & $n=8$ & $n=10$ & $n=12$ & $n=14$ & $n=16$ \\
    \midrule
    0.1 & 75\,\% & 90\,\% & 90\,\% & 80\,\% & 65\,\% & 65\,\% \\
    0.5 & 100\,\% & 100\,\% & 100\,\% & 100\,\% & 95\,\% & 100\,\% \\
    1, 2, 5 & 100\,\% & 100\,\% & 100\,\% & 100\,\% & 100\,\% & 100\,\% \\
    \bottomrule
  \end{tabular}
\end{table}

We use $\lambda = 0.5$ for all experiments.
At $\lambda = 0.1$ the optimiser frequently trades a conflict for better balance, especially at larger $n$, where more conflicts are available to trade.
Enforcing conflicts has a modest cost: at $\lambda = 0.5$ the optimal makespan is on average 4--9\,\% longer than the best makespan with conflicts ignored, depending on $n$.
% TODO(team): the sweep stops at n = 16, but E1 runs to n = 20 and E5 to n = 28.
% A spot check with the same rule found zero-violation optima on 9/10 instances at n = 20 and 3/4 at n = 24,
% and in results/baselines_timing.csv the exact (CP-SAT) optimum has 1 violation at n = 24 and n = 25.
% Either extend the sweep to n = 20, 24 and keep 0.5, or switch to lambda = 1 before the freeze.

\subsection{Worked example}

Consider three jobs with runtimes $\tilde w = (1, 1, 2)$, so $w = (0.5, 0.5, 1)$, one conflict $E = \{(0,1)\}$, and $\lambda = 1$ (chosen here for round numbers).
Eq.~\eqref{eq:ising} gives $J_{01} = 2(0.25) + 1 = 1.5$ and $J_{02} = J_{12} = 2(0.5)(1) = 1$, so $s = 1.5$ and
\[
  H_C = I + Z_0 Z_1 + \tfrac{2}{3} Z_0 Z_2 + \tfrac{2}{3} Z_1 Z_2 .
\]
Table~\ref{tab:example} lists all eight schedules.

\begin{table}[h]
  \centering
  \caption{All schedules for the worked example. Bit $i$ is job $i$; 0 = node A, 1 = node B. Each row is a mirror pair: the node columns describe the first bitstring, and the second swaps A and B.}
  \label{tab:example}
  \begin{tabular}{@{}cccccc@{}}
    \toprule
    Bits $x_0x_1x_2$ & Node A & Node B & $C(z)$ & Violations & Makespan \\
    \midrule
    000, 111 & all jobs & -- & 5 & 1 & 4 \\
    001, 110 & 0, 1 & 2 & 1 & 1 & 2 \\
    010, 101 & 0, 2 & 1 & \textbf{0} & 0 & 3 \\
    011, 100 & 0 & 1, 2 & \textbf{0} & 0 & 3 \\
    \bottomrule
  \end{tabular}
\end{table}

Without the conflict, the best schedule would put jobs 0 and 1 together against job 2, giving a perfect split with makespan 2.
The conflict forbids this, and the cost function instead selects the schedules with makespan 3 that keep jobs 0 and 1 apart.
The example also shows both effects from Section~\ref{sec:symmetry}: each schedule appears with its mirror image, and two different pairs of schedules share the minimum, so there are four optimal bitstrings.
